% Chapter 1: Principles & Mental Model
\chapter{Principles and Mental Model}
\label{ch:principles}

\begin{itshape}
You have just installed Claude Code. Before your first prompt, take five minutes
to understand what you are working with---not a chatbot, not an autocomplete engine,
but an agent with a specific architecture and specific constraints.
The principles in this chapter will inform every decision in this handbook.
\end{itshape}

\section{What Claude Code Actually Is}
\label{sec:mental-model}

\marginnote[Official]{Claude Code is an agentic coding tool that operates directly
in your terminal.\sourceurl{https://code.claude.com/docs/en/overview}}

Claude Code runs an \textbf{agent loop}\index{agent loop}: it reads your prompt, reasons about
the task, selects a tool (read a file, edit code, run a command), observes
the result, and decides what to do next. This cycle repeats until the task
is complete or Claude asks for clarification.

Three properties define the system:

\begin{description}
  \item[Context window]\index{context window} Everything Claude knows about your task---your prompt,
    files it has read, tool results, conversation history---occupies a finite
    window. When the window fills, older information fades. Context is not
    infinite memory; it is a workspace.

  \item[Tool use]\index{tool use} Claude does not type code into a file. It calls tools:
    \code{Read}, \code{Edit}, \code{Write}, \code{Bash}, \code{Glob},
    \code{Grep}. Each tool call consumes context (the request and the result)
    and produces observable effects. Understanding this matters because
    every tool call has a cost in context tokens.

  \item[Configuration layers] Claude's behavior is shaped by four layers:
    \code{CLAUDE.md} (project knowledge), rules (conditional context),
    settings (permissions and behavior), and output styles (communication patterns).
    These layers are not suggestions---they are the operating system
    through which Claude interprets your project.
\end{description}

\begin{whybox}[Why the Mental Model Matters]
Developers who treat Claude as a chatbot prompt it like one:
vague requests, no verification, no context management.
Developers who understand the agent loop prompt it like a system:
clear inputs, verification criteria, deliberate context control.
The same tool produces dramatically different results depending
on which mental model the developer holds.
\end{whybox}

% -------------------------------------------------------------------
\section{Four Engineering Principles}
\label{sec:four-principles}

These four principles underlie every recommendation in this handbook.
They apply to code written by AI, code written by humans,
and code written by both together.

\subsection{Never Fail Silently}

Every error must be explicitly reported with recovery options.\index{principles!never fail silently}
Silent failures are the most expensive kind of bug:
they produce incorrect results that look correct,
propagate through downstream systems undetected,
and surface only when the damage is difficult to reverse.

\begin{fullwidth}
\begin{beforeafter}
\before
\begin{minted}{markdown}
# CLAUDE.md
Handle errors gracefully.
\end{minted}

\after
\begin{minted}{markdown}
# CLAUDE.md
If you encounter an error, report it immediately
with the error message, your analysis of the root cause,
and 2-3 options for resolution.
Never work around errors silently.
\end{minted}
\end{beforeafter}
\end{fullwidth}

\margintip{The ``after'' version is longer but produces better outcomes.
Precision in error handling instructions compounds across every session.}

In AI-assisted development, this principle has an additional dimension:
Claude must be configured to surface problems rather than work around them.
A vague instruction to ``handle errors'' invites Claude to catch and suppress.
An explicit instruction to ``report errors with analysis'' makes every failure
visible and actionable.

\subsection{Simplicity Over Complexity}

Short functions. Flat structure. Self-documenting code.\index{principles!simplicity}
A 20-line function with a clear name is superior to a 5-line function
with three levels of abstraction.

\margintip{In notebooks, this means: one cell, one purpose.
Cells that do loading, cleaning, AND analysis violate this principle.}

\begin{whybox}[Why Simplicity Matters More with AI]
Simple code is easier for Claude to reason about, test, and modify correctly.
Complex abstractions increase the probability of subtle bugs in AI-generated
code because each layer of indirection is another opportunity for
misunderstanding. When Claude reads your codebase to learn conventions,
simple patterns propagate accurately; complex patterns propagate errors.
\end{whybox}

\subsection{Immutability by Default}

Return new data structures; mark mutations explicitly.\index{principles!immutability}
Pure functions are easier to test, easier to parallelize,
and easier for Claude to reason about.

When mutation is necessary (performance, I/O, state management),
make it explicit: name the function \code{update\_}, use mutable types,
and document the side effects. The reader---human or AI---should
never be surprised by what a function modifies.

\subsection{Fail Fast with Diagnostics}

Stop immediately on problems with full context:\index{principles!fail fast}
what failed, what was expected, what was received,
and what the caller can do about it.

\begin{minted}{python}
def process_data(
    df: pd.DataFrame, min_rows: int
) -> pd.DataFrame:
    if len(df) < min_rows:
        raise ValueError(
            f"Need {min_rows} rows, got {len(df)}. "
            f"Check data source or reduce "
            f"min_rows parameter."
        )
    # ... process
\end{minted}

\margintip{Error messages that include both the violation and a suggested fix
reduce debugging time dramatically.}

The error message includes: what went wrong (\code{Need N, got M}),
where to look (\code{data source}), and what to do about it
(\code{reduce min\_rows}). This is the difference between a
10-second fix and a 30-minute investigation.

% -------------------------------------------------------------------
\section{The Codebase Is the Curriculum}
\label{sec:codebase-curriculum}

\begin{keyinsight}
These principles are not Claude-specific---they are engineering principles
that become \emph{more} important when AI is generating code.
AI amplifies both good and bad patterns. If your codebase follows
these principles, Claude's contributions will follow them too.
If your codebase is silent about errors, full of mutable state,
and deeply abstracted, Claude's contributions will be the same.
The codebase is the curriculum. Make it teach the right lessons.
For data scientists, this includes your notebooks, pipeline scripts,
and feature engineering code---Claude learns patterns from all of them.
\end{keyinsight}

This insight has a practical consequence: the single most valuable
thing you can do before using Claude on an existing project is
to ensure that the code Claude will read exemplifies the patterns
you want Claude to produce. Configuration (\cref{ch:first-claude-md})
reinforces this, but the codebase itself is the primary teacher.

\marginnote[Cross-Ref]{For concrete techniques to make your codebase
a better teacher, see \cref{ch:thinking-partner}.}

% -------------------------------------------------------------------
\section{Quick Reference}

\begin{itemize}
  \item Claude Code is an agent loop: prompt $\rightarrow$ reason $\rightarrow$ tool $\rightarrow$ observe $\rightarrow$ repeat
  \item Context window is finite---every tool call costs tokens
  \item Four principles: never fail silently, simplicity, immutability, fail fast
  \item The codebase teaches Claude patterns---good and bad
  \item Configuration (next chapter) reinforces but does not replace good code
\end{itemize}

% -------------------------------------------------------------------
\begin{trybox}
Pick one function in your feature pipeline or data loading code---ideally
one Claude will read often. Audit it against the four principles:
\begin{enumerate}
  \item Does it fail silently anywhere? (e.g., silently dropping NaN rows
    without logging)
  \item Is it under 30 lines? If not, can you extract the transformation
    logic from the I/O?
  \item Does it mutate its input DataFrame? If so, is that explicit
    in the name?
  \item Do its error messages include what went wrong AND what to do about it?
\end{enumerate}
Fix one issue. This is your first investment in the codebase-as-curriculum.
\end{trybox}
